\section{Assignment}

\begin{enumerate}
	\item \textit{What are the four-strokes of a four-stroke petrol engine, and describe very briefly what happens during each stroke?}

	      \begin{enumerate}
		      \item Intake Stroke (TDC $\rightarrow$ BDC): At the Intake or Suction Stroke the inlet valve opens while the exhaust valve stays closed. As the piston moves downward, it creates a partial vacuum that draws a fresh air-fuel mixture into the cylinder.
		      \item Compression Stroke (BDC $\rightarrow$ TDC): At the Compression Stroke both valves close tightly to seal the cylinder. As the piston moves upward, it compresses the trapped air-fuel mixture into a clearance volume, raising it pressure and temperature.
		      \item Power Stroke (TDC $\rightarrow$ BDC): At the Power or Expansion Stroke the both valves remain still closed. The spark plug fires just before the TDC, igniting the compressed mixture. The rapid expansion of hot burning gases forces the piston downward, delivering mechanical work to the crankshaft.
		      \item Exhaust Stroke (BDC $\rightarrow$ TDC): At the Exhaust Stroke the exhaust valve opens while the inlet valve remains closed. As the piston moves upward again, it pushes the spent combustion gases out of the cylinder into the exhaust system.
	      \end{enumerate}

	\item \textit{Compare four-stroke engine to a two-stroke engine.}

	      \begin{table}[ht]
		      \centering
		      \begin{tabularx}{\textwidth}{X*{5}{>{\arraybackslash}X}}
			      \hline
			      \textbf{Four-Stroke}                                     & \textbf{Two-Stroke}                          \\
			      \hline
			      Four-Stroke (Intake, Compression, Power, Exhaust Stroke) & Two-Stroke (Upward \& Downward Strokes)      \\
			      Lower (4:1)                                              & Higher (1:1)                                 \\
			      Two revolutions per power stroke (720\degree)            & One revolution per power stroke (360\degree) \\
			      Complex                                                  & Simple                                       \\
			      High                                                     & Low                                          \\
			      Cleaner                                                  & Higher                                       \\
			      Longer engine life                                       & Shorter engine life                          \\
			      \hline
		      \end{tabularx}
		      \caption{Comparison of Four-Stroke and Two-Stroke Engines}
	      \end{table}

	\item \textit{Sketch a typical petrol engine and label all the components.}

	      The sketch of a typical petrol engine has been provided in Figure 02.

	\item \textit{Which engine produces more power? Explain.}

	      Generally, the two-stroke engines produce more power than the four-stroke engines. The reason is four-stroke engines produce one power stroke per cycle, while two-stroke engines produce two power strokes per cycle, providing more power to the crankshaft.

	\item \textit{Instead of valves, what is used in two-stroke engines?}

	      In two-stroke engines, instead of valves, ports are used.

	      \begin{enumerate}
		      \item Inlet Port: Inlet Port allows air-fuel mixture to enter the crankcase.
		      \item Outlet Port: Outlet Port transfers the air-fuel mixture to the cylinder.
		      \item Exhaust Port: Exhaust Port transfers the exhaust gases to the atmosphere.
	      \end{enumerate}

	\item \textit{What are the types of piston rings an how many rings do we use? How can piston rings help to optimize the piston clearance?}

	      There are two types of piston rings: compression and oil-controlling rings.

	      \begin{enumerate}
		      \item Compression Ring: Compression Ring reduces any leakage of gas between the piston and the cylinder.
		      \item Oil-Controlling Ring: Oil-Controlling Ring ensures excess oil is removed from the cylinder wall.
	      \end{enumerate}

	      In a engine, normally there are three rings, two compression rings and one oil-controlling ring.

	\item \textit{Briefly describe the Otto cycle and relate it to the petrol engine.}

	      The Otto cycle is an ideal thermodynamic cycle that describes how a spark-ignition engine converts chemical energy from fuel into mechanical energy. On a pressure-volume diagram, the ideal cycle has four main processes:

	      \begin{enumerate}
		      \item Adiabatic Compression (1 $\rightarrow$ 2)
		      \item Constant Volume Heat Addition (2 $\rightarrow$ 3)
		      \item Adiabatic Expansion (3 $\rightarrow$ 4)
		      \item Constant Volume Heat Removal (4 $\rightarrow$ 1)
	      \end{enumerate}

	      Relation to a Petrol Engine:

	      \begin{table}[h]
		      \centering
		      \begin{tabularx}{\textwidth}{X*{5}{>{\arraybackslash}X}}
			      \hline
			      \textbf{Otto Cycle} & \textbf{Petrol Engine} & \textbf{What Happen}                                         \\
			      \hline
			      Intake              & Suction                & Inlet valve opens, air and petrol vapour enter the cylinder. \\
			      Compression         & Compression            & Both valves close, the piston compresses the mixture.        \\

			      \hline
		      \end{tabularx}
	      \end{table}

	\item \textit{How does knocking occur in a Petrol/Diesel engine?}

	\item \textit{What are the measures taken in modern two-stroke petrol engines to reduce emissions and improve fuel efficiency?}
\end{enumerate}
